\subsection{Definition of formal power series. Order}

Let I be a set. We recall (III, p.~454 and 456) that the total algebra of the monoid \(\mathbf{N}^{(I)}\) over A is called the algebra of formal power series with respect to the indeterminates \(\mathbf{X}_i\) ( \(i \in I\) ) (or in the indeterminates \(\mathbf{X}_i\) ) with coefficients in A. It is denoted by \(\mathbf{A}[[\mathbf{X}_i]]_{i \in I}\) or \(\mathbf{A}[[(\mathbf{X}_i)_{i \in I}]]\) or also \(\mathbf{A}[[\mathbf{X}]]\) , on denoting by \(\mathbf{X}\) the family \((\mathbf{X}_i)_{i \in I}\) : in this paragraph we shall mainly use the notation \(\mathbf{A}[[I]]\) . Sometimes it is convenient to designate the canonical image in \(\mathbf{A}[[I]]\) of the element \(i\) of I by a symbol other than \(\mathbf{X}_i\) , for example \(\mathbf{Y}_i\) , \(\mathbf{Z}_i\) , \(\mathrm{T}_i\) , \ldots; the conventions used in this case are analogous to those for polynomials (IV, p.~1). The algebra \(\mathbf{A}[[I]]\) is then designated by \(\mathbf{A}[[\mathbf{Y}_i]]_{i \in I}\) , or \(\mathbf{A}[[\mathbf{Y}]]\) etc.

When I is a finite set of p elements, we also say that A{[}{[}I{]}{]} is an algebra of formal power series in p indeterminates. These algebras are all isomorphic, for fixed p.~An algebra of formal power series in 1, 2, \ldots{} indeterminates will also be denoted by A{[}{[}X{]}{]}, A{[}{[}U, V{]}{]}, \ldots, the set I of indices not being specified.

A formal power series \(u\) is conventionally written \(u = \sum_{\nu \in \mathbb{N}^{(I)}} \alpha_{\nu} X^{\nu}\) (cf.~IV, p.~1).

The \(\alpha_{\nu}\) are the coefficients of \(u\) ; there may be infinitely many of them \(\neq 0\) . The \(\alpha_{\nu}X^{\nu}\) are called terms of \(u\) ; for \(u\) to be a polynomial it is necessary and sufficient that \(u\) should have only a finite number of terms \(\neq 0\) . The terms \(\alpha_{\nu}X^{\nu}\) such that \(|\nu|=p\) are called the terms of total degree p.~The formal power series \(u_{p}=\sum_{|\nu|=p}\alpha_{\nu}X^{\nu}\) is called the homogeneous component of degree p of u (it is a polynomial when I is finite); \(u_{0}\) is identified with an element of A called also the constant term of u. We say that u is homogeneous of degree p if \(u = u_{p}\) . If u, \(v \in A[[I]]\) and w = uv, we have

\[
w _ {p} = \sum_ {q + r = p} u _ {q} v _ {r}\tag{1}
\]

for every integer \(p \geqslant 0\) .

We recall (III, p.~456), that the order \(\omega(u)\) of a formal power series \(u \neq 0\) is the least integer \(p\) such that \(u_p \neq 0\) . We shall agree to adjoin to \(\mathbf{Z}\) an element written \(\infty\) and extend the order relation and addition from \(\mathbf{Z}\) to \(\mathbf{Z} \cup \{\infty\}\) by the conventions

\[
n <   \infty , \quad \infty + \infty = \infty , \quad \infty + n = n + \infty = \infty
\]

for every \(n \in \mathbf{Z}\) ; we also put \(\omega(0) = \infty\) . With these conventions we have the relations

\[
\begin{array}{c} \omega (u + v) \geqslant \inf (\omega (u), \omega (v)), \\ \omega (u + v) = \inf (\omega (u), \omega (v)) \quad \text { if } \quad \omega (u) \neq \omega (v), \\ \omega (u v) \geqslant \omega (u) + \omega (v), \end{array}
\]

for any formal power series \(u\) and \(v\) in A{[}{[}I{]}{]}.

We recall (III, p.~457) that for any subset \(J\) of \(I\) we identify \(A[[I]]\) with \(A[[I - J]][[J]]\) , which allows us to define the order \(\omega_{j}(u)\) of a formal power series with respect to the \(X_{j}\) ( \(j \in J\) ), the homogeneous component of \(u\) with respect to the \(X_{j}\) ( \(j \in J\) ) etc.

Let \(\varphi\) be a homomorphism of A into a ring B. We extend \(\varphi\) to a homomorphism \(\overline{\varphi}\) of A{[}{[}I{]}{]} into B{[}{[}I{]}{]} by letting each formal power series \(u = \sum_{\nu} \alpha_{\nu} X^{\nu}\) correspond

to the formal power series \(\sum_{\nu} \varphi(\alpha_{\nu}) X^{\nu}\) ; we say that the latter is obtained by

applying \(\varphi\) to the coefficients of the formal power series \(u\) . We shall sometimes write \({}^{\varphi}u\) for \(\bar{\varphi}(u)\) .

In particular if A is a subring of B and \(\varphi\) the canonical injection of A into B, then the homomorphism \(\bar{\varphi}\) of A{[}{[}I{]}{]} into B{[}{[}I{]}{]} is injective; we shall in general identify A{[}{[}I{]}{]} with a subring of B{[}{[}I{]}{]} by means of \(\bar{\varphi}\) .

